\chapter{The Hodge Conjecture}\label{hodge-conjecture}
\chapterstatus{complete}

\section{Introduction}\label{hodge-conjecture:section-introduction}

On a nonsingular complex projective variety there are two ways of producing rational cohomology classes of type $(p,p)$. The algebro-geometric way takes
the class of a subvariety of codimension $p$ (Section~\ref{intersection-theory:section-cycle-class}). The Hodge-theoretic way takes all rational classes whose type is $(p,p)$. The first
set is contained in the second (Proposition~\ref{hodge-conjecture:proposition-easy-inclusion}). The Hodge conjecture asserts that, after taking $\QQ$-linear
combinations, they are equal. The inclusion is easy; the conjecture is a \emph{construction problem}: from a class satisfying two linear conditions one must produce
subvarieties, and there is no general method for producing subvarieties of codimension $\geq 2$. This chapter states the conjecture and its equivalent forms, proves the
elementary reductions, tells the history of its two corrections, and records the reasons for believing it and the reasons it is hard. We follow Deligne's official
problem description \cite{deligne-clay} and the original papers cited below.

\noindent\textbf{Prerequisites.} Chapter~\ref{hodge-structures} (Hodge Structures) and Section~\ref{intersection-theory:section-cycle-class} (The Cycle Class Map).

Throughout, $X$ is a nonsingular complex projective variety of dimension $n$, identified with the compact K\"ahler manifold $X^{\mathrm{an}}$ (Theorem~\ref{gaga:theorem-projective-kahler}),
and $0 \leq p \leq n$.

\section{Algebraic classes and Hodge classes}\label{hodge-conjecture:section-classes}

\begin{definition}\label{hodge-conjecture:definition-algebraic-classes}
The group of \emph{rational Hodge classes} of codimension $p$ is $\mathrm{Hdg}^p(X) = H^{2p}(X; \QQ) \cap H^{p,p}(X)$ (Definition~\ref{intersection-theory:definition-hodge-classes}). The space
of \emph{algebraic classes} is the image of the cycle class map,
\[
Z^p_{\mathrm{alg}}(X) = \mathrm{cl}(\CH^p(X)) \otimes \QQ = \Big\{\sum_ia_i\,\mathrm{cl}(Z_i) : a_i \in \QQ,\ Z_i \subseteq X \text{ subvarieties of codimension } p\Big\} \subseteq H^{2p}(X; \QQ).
\]
\end{definition}

\begin{proposition}[The easy inclusion]\label{hodge-conjecture:proposition-easy-inclusion}
$Z^p_{\mathrm{alg}}(X) \subseteq \mathrm{Hdg}^p(X)$.
\end{proposition}

\begin{proof}
This is Corollary~\ref{intersection-theory:corollary-hodge-classes}. The content is that a form of type $(a, b)$ with $a > \dim Z$ or $b > \dim Z$ restricts to zero on $Z_{\mathrm{reg}}$
(Proposition~\ref{kahler:proposition-cycle-classes-type}).
\end{proof}

\begin{conjecture}[Hodge conjecture]\label{hodge-conjecture:conjecture-hodge}
For every nonsingular complex projective variety $X$ and every $p$,
\[
Z^p_{\mathrm{alg}}(X) = \mathrm{Hdg}^p(X).
\]
That is, every rational cohomology class of type $(p,p)$ on $X$ is a rational linear combination of classes of algebraic subvarieties of codimension $p$. We write $\mathrm{HC}(X, p)$
for the statement for given $X$ and $p$, and $\mathrm{HC}(X)$ for all $p$.
\end{conjecture}

\begin{remark}\label{hodge-conjecture:remark-load-bearing}
Every hypothesis is necessary. The coefficients must be rational: the integral version is false, both for torsion classes (Atiyah--Hirzebruch) and for non-torsion classes (Koll\'ar).
$X$ must be projective: the analogous statement for compact K\"ahler manifolds is false (Voisin). Hodge's original, more general conjecture on the whole cohomology was false as stated
(Grothendieck). These failures are the subject of Chapter~\ref{known-cases} and Section~\ref{hodge-conjecture:section-history}.
\end{remark}

\begin{proposition}[Equivalent forms]\label{hodge-conjecture:proposition-equivalent-forms}
For nonsingular projective $X$ the following are equivalent to $\mathrm{HC}(X, p)$:
\begin{enumerate}
\item (homological form) every class in $H_{2n-2p}(X; \QQ)$ whose Poincar\'e dual is of type $(p,p)$ is a rational combination of fundamental classes of subvarieties of dimension $n - p$;
\item (Hodge-theoretic form) the map $\CH^p(X)_\QQ \to \Hom_{\mathrm{HS}}(\QQ(0), H^{2p}(X; \QQ)(p))$ induced by $\mathrm{cl}$ is surjective;
\item (vector bundle form) $\mathrm{Hdg}^p(X)$ is contained in the $\QQ$-algebra generated by the Chern classes of algebraic vector bundles on $X$.
\end{enumerate}
\end{proposition}

\begin{proof}
(1) is Poincar\'e duality, which maps fundamental classes to cycle classes (Definition~\ref{intersection-theory:definition-cycle-class}) and Hodge classes to the corresponding classes in
homology. (2) holds because $\Hom_{\mathrm{HS}}(\QQ(0), H^{2p}(p))$ is the space of rational classes of type $(0,0)$ in $H^{2p}(p)$, which is $\mathrm{Hdg}^p(X)$ twisted
(Definition~\ref{hodge-structures:definition-hodge-classes}). (3) follows from Proposition~\ref{intersection-theory:proposition-chern-character-cycle}: the $\QQ$-span of algebraic classes of all
codimensions equals the $\QQ$-algebra generated by Chern classes of vector bundles, and its degree-$2p$ part is $Z^p_{\mathrm{alg}}(X)$.
\end{proof}

\begin{remark}\label{hodge-conjecture:remark-correspondence-form}
Applied to a product $X \times Y$, the conjecture predicts that every morphism of Hodge structures $H^k(X; \QQ) \to H^{k+2r-2n}(Y; \QQ)(r - n)$ that is a Hodge class in
$H^{2r}(X \times Y; \QQ)$ (via K\"unneth and Poincar\'e duality) is induced by an algebraic correspondence (Definition~\ref{intersection-theory:definition-correspondence}). This is the form in which
the conjecture is used: for instance, it would make the K\"unneth projectors and the Lefschetz inverse algebraic (Chapter~\ref{absolute-hodge}).
\end{remark}

\section{First reductions and elementary cases}\label{hodge-conjecture:section-reductions}

\begin{proposition}\label{hodge-conjecture:proposition-trivial-degrees}
$\mathrm{HC}(X, 0)$ and $\mathrm{HC}(X, n)$ hold: $\mathrm{Hdg}^0(X) = H^0(X; \QQ)$ is spanned by the classes of the connected components, and $\mathrm{Hdg}^n(X) = H^{2n}(X; \QQ)$ is spanned by the classes of
points, one in each connected component.
\end{proposition}

\begin{proof}
$H^0$ and $H^{2n}$ are entirely of type $(0,0)$ and $(n,n)$. The class of a connected component is its fundamental class, and the class of a point generates $H^{2n}$ of its component
(Example~\ref{intersection-theory:example-classes}(3)).
\end{proof}

\begin{proposition}[Hard Lefschetz reduction]\label{hodge-conjecture:proposition-lefschetz-reduction}
Let $2p \leq n$. Then $\mathrm{HC}(X, p)$ implies $\mathrm{HC}(X, n - p)$.
\end{proposition}

\begin{proof}
Let $h = \mathrm{cl}(H)$ be the class of a hyperplane section, an algebraic class. By Theorem~\ref{kahler:theorem-hard-lefschetz}, $L^{n-2p} \colon H^{2p}(X; \QQ) \to H^{2n-2p}(X; \QQ)$,
$\alpha \mapsto h^{n-2p} \cup \alpha$, is an isomorphism; it is defined over $\QQ$ and of type $(n-2p, n-2p)$, so it maps $\mathrm{Hdg}^p(X)$ isomorphically onto $\mathrm{Hdg}^{n-p}(X)$ (its inverse is
also of type $(-(n-2p), -(n-2p))$ and defined over $\QQ$). If $\beta \in \mathrm{Hdg}^{n-p}(X)$, then $\beta = h^{n-2p}\alpha$ with $\alpha \in \mathrm{Hdg}^p(X)$, which is algebraic by hypothesis, and
$\beta$ is a product of algebraic classes, hence algebraic (Theorem~\ref{intersection-theory:theorem-cycle-class}(1)).
\end{proof}

\begin{remark}\label{hodge-conjecture:remark-small-dimension}
The case $p = 1$ is the Lefschetz $(1,1)$-theorem (Section~\ref{hodge-conjecture:section-lefschetz-11}). With Proposition~\ref{hodge-conjecture:proposition-lefschetz-reduction} it gives $\mathrm{HC}(X, n - 1)$, and
therefore the Hodge conjecture for all varieties of dimension at most $3$. The first open case is $p = 2$ on fourfolds. The converse of
Proposition~\ref{hodge-conjecture:proposition-lefschetz-reduction}, that $\mathrm{HC}(X, n-p)$ implies $\mathrm{HC}(X, p)$, is not known in general: it would require the inverse of $L^{n-2p}$ to be
induced by an algebraic correspondence, which is the Lefschetz standard conjecture (Chapter~\ref{absolute-hodge}).
\end{remark}

\begin{proposition}\label{hodge-conjecture:proposition-dominant-maps}
Let $f \colon X \to Y$ be a surjective morphism of nonsingular projective varieties of the same dimension. If $\mathrm{HC}(X, p)$ holds, then $\mathrm{HC}(Y, p)$ holds.
\end{proposition}

\begin{proof}
$f$ is generically finite of some degree $d \geq 1$, so $f_*[X] = d[Y]$ in $\CH^0(Y)$ (Theorem~\ref{intersection-theory:theorem-functoriality}). The pullback $f^*$ and the Gysin pushforward $f_*$ in
cohomology are morphisms of Hodge structures (of type $(0,0)$ between the relevant degrees, as $\dim X = \dim Y$), and both map algebraic classes to algebraic classes
(Theorem~\ref{intersection-theory:theorem-cycle-class}(2)). By the projection formula, $f_*f^*\beta = f_*(1) \cup \beta = d\beta$. So for $\beta \in \mathrm{Hdg}^p(Y)$, $f^*\beta \in \mathrm{Hdg}^p(X)$ is algebraic,
and $\beta = \frac1df_*(f^*\beta)$ is algebraic.
\end{proof}

\begin{example}\label{hodge-conjecture:example-cellular}
If $X$ has a filtration by closed subvarieties $X = X_0 \supseteq X_1 \supseteq \cdots \supseteq X_m = \emptyset$ with each $X_i \setminus X_{i+1}$ a disjoint union of affine spaces, then
$H^{\mathrm{odd}}(X; \ZZ) = 0$ and $H^{2p}(X; \ZZ)$ is free on the classes of the closures of the cells of codimension $p$ \cite[Example 19.1.11]{fulton-intersection}. So $\mathrm{HC}(X)$
holds, even integrally. This covers projective spaces, Grassmannians, flag varieties and smooth projective toric varieties; here the Hodge structures are of Tate type and there is
no room for a non-algebraic Hodge class. For $\PP^n$ it is Example~\ref{intersection-theory:example-classes}(1).
\end{example}

\begin{example}\label{hodge-conjecture:example-elliptic-product}
Let $E$, $E'$ be complex elliptic curves. By K\"unneth,
$H^2(E \times E'; \QQ) = H^2(E) \oplus (H^1(E) \otimes H^1(E')) \oplus H^2(E')$, and the outer summands are spanned by the classes of $\{\mathrm{pt}\} \times E'$ and $E \times \{\mathrm{pt}\}$.
The Hodge classes in the middle summand correspond to morphisms of Hodge structures $H_1(E) \to H_1(E')$ (after Poincar\'e duality on $E$), that is, by
Theorem~\ref{hodge-structures:theorem-weight-one}, to elements of $\Hom(E, E') \otimes \QQ$. The Hodge classes are nonzero exactly when $E$ and $E'$ are isogenous. Each homomorphism
$\varphi$ has a graph $\Gamma_\varphi \subseteq E \times E'$, whose class has K\"unneth component in $H^1 \otimes H^1$ equal to the class corresponding to $\varphi$, so $\mathrm{HC}(E \times E', 1)$ holds.
(This is also a consequence of the Lefschetz $(1,1)$-theorem.)
\end{example}

\begin{remark}\label{hodge-conjecture:remark-products}
$\mathrm{HC}(X)$ and $\mathrm{HC}(Y)$ do not formally imply $\mathrm{HC}(X \times Y)$. The K\"unneth summands $H^i(X) \otimes H^j(Y)$ with $i$ and $j$ odd contain Hodge classes that are not
products of Hodge classes on the factors. Example~\ref{hodge-conjecture:example-elliptic-product} is the simplest case, and the Weil classes on abelian fourfolds
(Chapter~\ref{known-cases}) are the serious ones. For blow-ups, on the other hand, the blow-up formula $H^\bullet(\tilde X) = H^\bullet(X) \oplus \bigoplus_{j=1}^{c-1}H^{\bullet-2j}(W)(-j)$ for the blow-up
along a nonsingular centre $W$ of codimension $c$ is induced by algebraic correspondences. So $\mathrm{HC}$ for $X$ and $W$ in the relevant degrees gives it for $\tilde X$ \cite[Theorem 7.31]{voisin-hodge-1}.
\end{remark}

\section{History and two corrections}\label{hodge-conjecture:section-history}

\begin{remark}[Lefschetz]\label{hodge-conjecture:remark-lefschetz}
In 1924 Lefschetz proved that on a nonsingular projective surface an integral class of type $(1,1)$ in $H^2$ is the class of a divisor, using normal functions. The modern proof via the
exponential sequence is due to Kodaira and Spencer \cite{kodaira-spencer-1953}. Lefschetz also proved the hyperplane theorem and, with a gap, the hard Lefschetz theorem, whose first complete
proof is Hodge's, by harmonic forms. The framework of the conjecture (hyperplane sections, primitive cohomology, divisors versus cohomology) is Lefschetz's.
\end{remark}

\begin{remark}[Hodge]\label{hodge-conjecture:remark-hodge-1950}
At the International Congress of 1950 Hodge proposed a generalisation of Lefschetz's theorem to all codimensions \cite{hodge-1950}. He stated it with integral coefficients, and he also
proposed a stronger conjecture about the whole cohomology. Both versions were false, for different reasons.
\end{remark}

\begin{remark}[First correction: integral to rational]\label{hodge-conjecture:remark-integral}
Atiyah and Hirzebruch showed in 1962 that there are nonsingular projective varieties with torsion classes in $H^{2p}(X; \ZZ)$ (automatically Hodge classes) that are not algebraic
\cite{atiyah-hirzebruch-1962}. The obstruction is topological: algebraic classes come from complex vector bundles, so the differentials of the Atiyah--Hirzebruch spectral sequence, starting
with $d_3 = \mathrm{Sq}^3_\ZZ$, must vanish on them. Koll\'ar later gave non-torsion counterexamples: on a very general hypersurface of suitable degree in $\PP^4$ the generator of
$H^4(X; \ZZ) \cong \ZZ$ is not algebraic, although a multiple of it is (Chapter~\ref{known-cases}). Rational coefficients remove both phenomena.
\end{remark}

\begin{definition}\label{hodge-conjecture:definition-coniveau}
The \emph{coniveau filtration} on $H^k(X; \QQ)$ is
\[
N^cH^k(X; \QQ) = \sum_Z\ker\big(H^k(X; \QQ) \to H^k(X \setminus Z; \QQ)\big),
\]
the sum over closed algebraic subsets $Z \subseteq X$ of codimension $\geq c$: the classes \emph{supported in codimension $c$}.
\end{definition}

\begin{proposition}\label{hodge-conjecture:proposition-coniveau}
$N^cH^k(X; \QQ)$ is a sub-Hodge structure of $H^k(X; \QQ)$ contained in $F^cH^k(X; \CC)$. Moreover $N^pH^{2p}(X; \QQ) = Z^p_{\mathrm{alg}}(X)$.
\end{proposition}

\begin{proof}[Sketch of proof]
A class supported on $Z$ lies in the image of the Gysin map $H^{k-2c'}(\tilde Z; \QQ)(-c') \to H^k(X; \QQ)$ from a resolution of singularities $\tilde Z$ of a component of $Z$ of codimension
$c' \geq c$, by Deligne's theory of weights; the Gysin map is a morphism of Hodge structures of type $(c', c')$, so its image
is a sub-Hodge structure contained in $F^{c'} \subseteq F^c$ \cite{deligne-hodge-2}. For $k = 2p$ and $c = p$, the Gysin image of $H^0(\tilde Z)$ is spanned by $\mathrm{cl}(Z)$.
\end{proof}

\begin{remark}[Second correction: Grothendieck's amendment]\label{hodge-conjecture:remark-grothendieck}
Hodge's general conjecture asserted $N^cH^k(X; \QQ) = H^k(X; \QQ) \cap F^cH^k(X; \CC)$. In 1969 Grothendieck observed that this is false for trivial reasons \cite{grothendieck-1969}: the left side is a
sub-Hodge structure, while the right side, a rational subspace contained in $H^{k-c,c} \oplus \cdots \oplus H^{c,k-c}$, need not be one, and he gave an abelian variety where they differ. He
proposed the \emph{generalised Hodge conjecture}: $N^cH^k(X; \QQ)$ is the largest sub-Hodge structure of $H^k(X; \QQ)$ contained in $F^cH^k(X; \CC)$. For $k = 2p$ and $c = p$ the largest
sub-Hodge structure contained in $F^p \cap \overline{F^p} = H^{p,p}$ is spanned by $\mathrm{Hdg}^p(X)$, so the generalised conjecture contains Conjecture~\ref{hodge-conjecture:conjecture-hodge}.
\end{remark}

\begin{remark}[The Clay problem]\label{hodge-conjecture:remark-clay}
In 2000 the Hodge conjecture became one of the seven Millennium Prize Problems of the Clay Mathematics Institute, with an official description by Deligne \cite{deligne-clay}, which emphasises
the reformulation through absolute Hodge classes (Chapter~\ref{absolute-hodge}).
\end{remark}

\section{Why one might believe it, and why it is hard}\label{hodge-conjecture:section-evidence}

\begin{remark}[Evidence]\label{hodge-conjecture:remark-evidence}
\begin{enumerate}
\item \emph{It holds where it can be checked}: for $p \in \{0, 1, n-1, n\}$, hence in dimension $\leq 3$; for cellular varieties; for cubic fourfolds \cite{zucker-1977}; for many abelian
varieties (Chapter~\ref{known-cases}). No verified case has failed.
\item \emph{Hodge loci are algebraic}. In a family of projective manifolds, the locus where a given integral class stays of type $(p,p)$ is a countable union of algebraic subvarieties, by the
theorem of Cattani, Deligne and Kaplan \cite{cattani-deligne-kaplan}. This is what the conjecture predicts, since loci where cycles exist are algebraic (Hilbert schemes), and the theorem is
proved without it (Chapter~\ref{vhs}).
\item \emph{Hodge classes on abelian varieties are absolute} \cite{deligne-abelian}: they are compatible with all automorphisms of $\CC$, as algebraic classes are and a random Hodge class would not be
(Chapter~\ref{absolute-hodge}).
\item \emph{Coherence}: the conjecture follows from, and is required by, the expected properties of pure motives and the standard conjectures (Chapter~\ref{absolute-hodge}).
\end{enumerate}
\end{remark}

\begin{remark}[Difficulty]\label{hodge-conjecture:remark-difficulty}
\begin{enumerate}
\item \emph{No construction machine in codimension $\geq 2$.} For divisors, an integral class of type $(1,1)$ lifts to a line bundle by the exponential sequence, and a high twist of a line bundle
has sections whose zero loci are divisors (Section~\ref{hodge-conjecture:section-lefschetz-11}). This uses that $\Pic(X) = H^1(X, \mathcal{O}_X^*)$ is the cohomology of a sheaf of abelian groups. In codimension $p \geq 2$,
Bloch's formula $\CH^p(X) \cong H^p(X, \mathcal{K}_p)$ expresses the Chow group through a sheaf of $K$-groups, but this transfers the difficulty rather than removing it.
\item \emph{Deformation theory obstructs.} Hilbert schemes of subvarieties of codimension $\geq 2$ are typically singular of excess dimension, and deformations of a cycle do not see the target
Hodge class.
\item \emph{The transcendental input is weak.} Hodge theory says nothing about where in $X$ a cycle should lie. Griffiths' normal functions (Chapter~\ref{vhs}) convert the problem
into the existence of singularities of certain sections, which is as hard in general.
\item \emph{The conjecture is close to false.} Its integral, K\"ahler and general versions all fail (Chapter~\ref{known-cases}). A proof must distinguish it from these neighbours.
\end{enumerate}
\end{remark}

\section{Exercises}\label{hodge-conjecture:section-exercises}

\begin{exercise}\label{hodge-conjecture:exercise-odd-weight}
Show that a $\QQ$-Hodge structure of odd weight has no Hodge classes in any Tate twist, so that Hodge classes live only in even degrees. Explain why the K\"unneth summands
$H^{\mathrm{odd}}(X) \otimes H^{\mathrm{odd}}(Y)$ of a product are where new Hodge classes appear.
\end{exercise}

\begin{solution}
$V(m)$ has weight $k - 2m$, odd, so it has no component of type $(0,0)$, and $\Hom_{\mathrm{HS}}(\QQ(0), V(m)) = 0$. The tensor product of two odd-weight structures has even weight, and it can
contain Hodge classes not coming from Hodge classes on the factors, which have none in odd degree (Example~\ref{hodge-conjecture:example-elliptic-product}).
\end{solution}

\begin{exercise}\label{hodge-conjecture:exercise-projective-product}
Verify $\mathrm{HC}(\PP^m \times \PP^k)$ directly.
\end{exercise}

\begin{solution}
By K\"unneth, $H^\bullet = \QQ[a, b]/(a^{m+1}, b^{k+1})$ with $a$, $b$ the pullbacks of hyperplane classes, all classes are of Tate type and spanned by the monomials $a^ib^j$, which are the classes
of the products of linear subspaces $L \times L'$ (Example~\ref{intersection-theory:example-classes}(1) and Theorem~\ref{intersection-theory:theorem-cycle-class}(2)).
\end{solution}

\begin{exercise}\label{hodge-conjecture:exercise-finite-quotient}
Let a finite group $G$ act on $X$ with nonsingular projective quotient $Y = X/G$. Show that $\mathrm{HC}(X, p)$ implies $\mathrm{HC}(Y, p)$.
\end{exercise}

\begin{solution}
The quotient map $X \to Y$ is surjective and finite, of degree $|G|$, so Proposition~\ref{hodge-conjecture:proposition-dominant-maps} applies.
\end{solution}

\section{The Lefschetz (1,1) Theorem}\label{hodge-conjecture:section-lefschetz-11}
The Lefschetz $(1,1)$-theorem is the one case of the Hodge conjecture that is completely understood, and it holds even with integral coefficients: on a
nonsingular projective variety every integral class of type $(1,1)$ in $H^2$ is the class of a divisor (Theorem~\ref{hodge-conjecture:theorem-lefschetz-11}). Its proof
uses the exponential sequence, which exists because the Picard group is the cohomology of a sheaf of abelian groups. With hard Lefschetz it gives the Hodge conjecture
in codimension $n - 1$, hence in dimension at most three. The following sections prove the theorem, draws the consequences, and explains precisely which step has no analogue in
higher codimension. References: \cite[Chapter 1, Section 1]{griffiths-harris}, \cite[Section 11.3]{voisin-hodge-1}, \cite{kodaira-spencer-1953}.

\noindent\textbf{Prerequisites.} Chapter~\ref{hodge-conjecture} (The Hodge Conjecture) and Chapter~\ref{gaga} (Analytic Spaces and GAGA).

Throughout, $X$ is a compact K\"ahler manifold of dimension $n$; from Section~\ref{hodge-conjecture:section-projective} on it is projective.

\section{The theorem for compact K\"ahler manifolds}\label{hodge-conjecture:section-kahler}

\begin{theorem}[Lefschetz $(1,1)$, K\"ahler form]\label{hodge-conjecture:theorem-kahler}
Let $X$ be a compact K\"ahler manifold. The image of the first Chern class $c_1 \colon \Pic(X) \to H^2(X; \ZZ)$ is exactly the group of integral Hodge classes
$\mathrm{Hdg}^1(X; \ZZ)$, the classes whose image in $H^2(X; \CC)$ is of type $(1,1)$.
\end{theorem}

\begin{proof}
The Chern class of a holomorphic line bundle is represented by $\frac{i}{2\pi}F$, a closed real $(1,1)$-form (Theorem~\ref{complex-manifolds:theorem-chern-connection}), so
$c_1(\Pic(X)) \subseteq \mathrm{Hdg}^1(X; \ZZ)$. Conversely, let $a \in H^2(X; \ZZ)$ have image of type $(1,1)$. By the \emph{key lemma} (Lemma~\ref{kodaira:lemma-integral-11}), the map
$H^2(X; \ZZ) \to H^2(X, \mathcal{O}_X)$ induced by $\ZZ_X \subseteq \mathcal{O}_X$ is the composite of $H^2(X; \ZZ) \to H^2(X; \CC)$ with the projection to $H^{0,2}(X) \cong H^2(X, \mathcal{O}_X)$.
So $a$ maps to $0$ in $H^2(X, \mathcal{O}_X)$. In the long exact sequence of the exponential sequence $0 \to \ZZ \to \mathcal{O}_X \to \mathcal{O}_X^* \to 0$
(Theorem~\ref{complex-manifolds:theorem-exponential}),
\[
H^1(X, \mathcal{O}_X^*) = \Pic(X) \xrightarrow{c_1} H^2(X; \ZZ) \to H^2(X, \mathcal{O}_X),
\]
so $a$ is in the image of $c_1$.
\end{proof}

\begin{corollary}\label{hodge-conjecture:corollary-neron-severi}
$\mathrm{NS}(X) := c_1(\Pic(X)) = \mathrm{Hdg}^1(X; \ZZ)$, and $\Pic(X)$ is an extension $0 \to \Pic^0(X) \to \Pic(X) \to \mathrm{NS}(X) \to 0$ of a finitely generated abelian group by the complex torus
$\Pic^0(X) = H^1(X, \mathcal{O}_X)/H^1(X; \ZZ)$ of dimension $h^{0,1} = b_1/2$. The \emph{Picard number} $\rho(X) = \operatorname{rank}\mathrm{NS}(X)$ satisfies $\rho(X) \leq h^{1,1}(X)$.
\end{corollary}

\begin{proof}
The first statement is Theorem~\ref{hodge-conjecture:theorem-kahler}, and the extension is the exponential sequence. $\mathrm{NS}(X) \subseteq H^2(X; \ZZ)$ is finitely generated, and
$\Pic^0(X)$ is a torus by Remark~\ref{complex-manifolds:remark-picard-variety} and Exercise~\ref{kahler:exercise-h1-real}. Finally $\mathrm{NS}(X) \otimes \RR$ injects into the real
classes of type $(1,1)$, which form a real vector space of dimension $h^{1,1}$ (a real structure on $H^{1,1}$, since $\overline{H^{1,1}} = H^{1,1}$).
\end{proof}

\begin{proposition}\label{hodge-conjecture:proposition-flat}
Let $X$ be a compact K\"ahler manifold, and let $U(1)_X$ be the sheaf of locally constant functions with values in $U(1) = \{z \in \CC : |z| = 1\}$. The inclusion $U(1)_X \subseteq \mathcal{O}_X^*$ induces a
homomorphism $b \colon H^1(X, U(1)_X) \to \Pic(X)$, whose image consists of the \emph{flat unitary} line bundles, those given by locally constant transition functions of modulus one.
\begin{enumerate}
\item $b$ maps the kernel of the connecting map $\delta \colon H^1(X, U(1)_X) \to H^2(X; \ZZ)$ isomorphically onto $\Pic^0(X)$, and this kernel is $H^1(X; \RR)/H^1(X; \ZZ)$, a real torus of dimension $b_1(X)$.
\item A holomorphic line bundle is flat unitary if and only if its first Chern class is a torsion element of $H^2(X; \ZZ)$. In particular every torsion class in $H^2(X; \ZZ)$ lies in $\mathrm{NS}(X)$.
\end{enumerate}
\end{proposition}

\begin{proof}
The map $t \mapsto e^{2\pi it}$ gives a short exact sequence of sheaves $0 \to \ZZ_X \to \RR_X \to U(1)_X \to 0$: on a connected neighbourhood a locally constant function is constant, and $\RR \to U(1)$ is
surjective with kernel $\ZZ$. The inclusions $\RR_X \subseteq \mathcal{O}_X$ and $U(1)_X \subseteq \mathcal{O}_X^*$ map it to the exponential sequence of
Theorem~\ref{complex-manifolds:theorem-exponential}, and by naturality of the long exact sequence (Theorem~\ref{homological-algebra:theorem-long-exact-sequence}) we get a commutative diagram with exact rows
\begin{center}
\begin{tikzcd}
H^1(X; \ZZ) \arrow[r] \arrow[d, equal] & H^1(X; \RR) \arrow[r, "\beta"] \arrow[d, "a"] & H^1(X, U(1)_X) \arrow[r, "\delta"] \arrow[d, "b"] & H^2(X; \ZZ) \arrow[r] \arrow[d, equal] & H^2(X; \RR) \\
H^1(X; \ZZ) \arrow[r] & H^1(X, \mathcal{O}_X) \arrow[r, "\beta'"] & \Pic(X) \arrow[r, "c_1"] & H^2(X; \ZZ) &
\end{tikzcd}
\end{center}
where $a$ is induced by $\RR_X \subseteq \mathcal{O}_X$. By Exercise~\ref{kahler:exercise-h1-real}, $a$ is an isomorphism of real vector spaces. By exactness,
$\ker\delta = \beta(H^1(X; \RR)) \cong H^1(X; \RR)/H^1(X; \ZZ)$ and $\Pic^0(X) = \ker c_1 = \beta'(H^1(X, \mathcal{O}_X)) \cong H^1(X, \mathcal{O}_X)/H^1(X; \ZZ)$, and $a$ induces an isomorphism between these
quotients, since the left square commutes. As $b\beta = \beta'a$, this isomorphism is the restriction of $b$, proving (1); the image of $H^1(X; \ZZ)$ in $H^1(X; \RR)$ is a lattice by the universal coefficient
theorem (Theorem~\ref{singular-homology:theorem-universal-coefficients}).

(2) If $L = b(y)$, then $c_1(L) = \delta(y)$ maps to $0$ in $H^2(X; \RR)$ by exactness of the top row, so it is torsion, because the kernel of $H^2(X; \ZZ) \to H^2(X; \RR)$ is the torsion subgroup (again by
the universal coefficient theorem). Conversely, let $c_1(L)$ be torsion. Then it maps to $0$ in $H^2(X; \RR)$, so $c_1(L) = \delta(y)$ for some $y$, and $L \otimes b(y)^{-1}$ has $c_1 = 0$. By (1),
$L \otimes b(y)^{-1} = b(y')$ is flat unitary, and so is $L = b(yy')$. The same argument, starting from a torsion class $t = \delta(y)$, shows $t = c_1(b(y)) \in \mathrm{NS}(X)$.
\end{proof}

\begin{example}\label{hodge-conjecture:example-two-torus}
Let $T = \CC^2/\Lambda$ be a complex torus of dimension $2$ and $\sigma = dz_1 \wedge dz_2$. By Example~\ref{complex-manifolds:example-hodge-numbers}, $h^{2,0} = h^{0,2} = 1$ and $h^{1,1} = 4$, with $H^{2,0}(T)$
spanned by $[\sigma]$ and $H^{0,2}(T)$ by $[\bar\sigma]$; and $H^2(T; \ZZ) \cong \ZZ^6$ (Example~\ref{singular-homology:example-torus}). A real class $a \in H^2(T; \RR)$ is of type $(1,1)$ if and only if
$\int_Ta \wedge \sigma = 0$. Indeed, write $a = a^{2,0} + a^{1,1} + a^{0,2}$ (Theorem~\ref{kahler:theorem-hodge-decomposition}), with $a^{2,0} = \overline{a^{0,2}}$ because $a$ is real, and
$a^{0,2} = c[\bar\sigma]$. For type reasons only $a^{0,2}$ contributes to $\int_Ta \wedge \sigma$, so the integral is $c\int_T\bar\sigma \wedge \sigma$, and $\bar\sigma \wedge \sigma$ is a nonzero constant
multiple of $dx_1 \wedge dy_1 \wedge dx_2 \wedge dy_2$ ($z_j = x_j + iy_j$); hence it vanishes if and only if $c = 0$, that is, $a^{0,2} = a^{2,0} = 0$.

By Theorem~\ref{hodge-conjecture:theorem-kahler}, $\mathrm{NS}(T)$ is therefore the kernel of $\varphi \colon H^2(T; \ZZ) \to \CC$, $\varphi(a) = \int_Ta \wedge \sigma$. If $e_1, \ldots, e_6$ is a $\ZZ$-basis of
$H^2(T; \ZZ)$ with \emph{periods} $\pi_k = \varphi(e_k)$, then
\[
\rho(T) = 6 - \dim_\QQ(\QQ\pi_1 + \cdots + \QQ\pi_6),
\]
since the kernel of $\ZZ^6 \to \CC$, $m \mapsto \sum_km_k\pi_k$, has the same rank as the kernel of its extension to $\QQ^6$. The real extension $H^2(T; \RR) \to \CC$ has as kernel the real classes of type
$(1,1)$, of real dimension $h^{1,1} = 4$, so it is surjective and the $\pi_k$ span $\CC$ over $\RR$. Hence the $\QQ$-dimension above is at least $2$ and $\rho(T) \leq 4$, as predicted by
Corollary~\ref{hodge-conjecture:corollary-neron-severi}. The Picard number of $T$ is thus governed by the linear relations over $\QQ$ among six complex numbers: $\rho(T) = 4$ when the periods span only a
two-dimensional $\QQ$-vector space, as for $E \times E$ with $E = \CC/(\ZZ + \ZZ i)$ (Example~\ref{hodge-conjecture:example-neron-severi}(2), since $\Hom(E, E) \supseteq \ZZ[i]$ has rank $2$), and
$\rho(T) = 0$ when they are linearly independent over $\QQ$, which is the situation of Exercise~\ref{hodge-conjecture:exercise-picard-number-torus}.
\end{example}

\section{The projective case: classes of divisors}\label{hodge-conjecture:section-projective}

\begin{theorem}[Lefschetz $(1,1)$]\label{hodge-conjecture:theorem-lefschetz-11}
Let $X$ be a nonsingular complex projective variety. Every class in $\mathrm{Hdg}^1(X; \ZZ)$ is the class of a divisor: the cycle class map $\CH^1(X) \to \mathrm{Hdg}^1(X; \ZZ)$ is surjective. In particular
$\mathrm{HC}(X, 1)$ holds, even integrally.
\end{theorem}

\begin{proof}
By Theorem~\ref{hodge-conjecture:theorem-kahler}, a class $a \in \mathrm{Hdg}^1(X; \ZZ)$ is $c_1(L)$ for a holomorphic line bundle $L$ on $X^{\mathrm{an}}$. By GAGA
(Corollary~\ref{gaga:corollary-consequences}(2)), $L$ is the analytification of an algebraic line bundle, still denoted $L$. Let $H$ be a hyperplane section. By Serre's theorem
(Theorem~\ref{schemes:theorem-serre-generation}), $L(m) = L \otimes \mathcal{O}(mH)$ is globally generated for $m \gg 0$, so it has a nonzero section, whose divisor $D_1$ is effective with
$\mathcal{O}(D_1) \cong L(m)$ (Proposition~\ref{divisors:proposition-cartier-picard}). Likewise $\mathcal{O}(mH)$ has the divisor $D_2 = mH$. So $L \cong \mathcal{O}(D_1 - D_2)$, and
$a = c_1(L) = \mathrm{cl}(D_1) - \mathrm{cl}(D_2)$ by Example~\ref{intersection-theory:example-classes}(2).
\end{proof}

\begin{corollary}\label{hodge-conjecture:corollary-codimension-n-1}
For a nonsingular projective variety $X$ of dimension $n$, $\mathrm{HC}(X, n - 1)$ holds. Consequently the Hodge conjecture holds for all nonsingular projective varieties of
dimension at most $3$.
\end{corollary}

\begin{proof}
$\mathrm{HC}(X, 1)$ holds by Theorem~\ref{hodge-conjecture:theorem-lefschetz-11}, and implies $\mathrm{HC}(X, n - 1)$ for $n \geq 2$ by Proposition~\ref{hodge-conjecture:proposition-lefschetz-reduction}.
Together with the trivial cases $p = 0$ and $p = n$ (Proposition~\ref{hodge-conjecture:proposition-trivial-degrees}), this covers all $p$ when $n \leq 3$.
\end{proof}

\begin{remark}\label{hodge-conjecture:remark-integral-curves}
For curves on a threefold, the integral version fails in general: Proposition~\ref{hodge-conjecture:proposition-lefschetz-reduction} only gives that some multiple of an integral Hodge
class in $H^4$ is algebraic, since $L^{n-2}$ need not be an isomorphism over $\ZZ$. Koll\'ar's counterexamples to the integral Hodge conjecture are integral classes in $H^4$ of threefolds
(Chapter~\ref{known-cases}). On the other hand, Voisin proved the integral Hodge conjecture for curve classes on uniruled threefolds and on Calabi--Yau threefolds \cite{voisin-2006}.
\end{remark}

\begin{example}\label{hodge-conjecture:example-neron-severi}
\begin{enumerate}
\item For $X \subseteq \PP^3$ a nonsingular quartic surface, $\mathrm{NS}(X)$ is a sublattice of $H^2(X; \ZZ) \cong \ZZ^{22}$ of rank $\rho \in [1, 20]$, and $\rho = 20$ is attained (for instance by the Fermat quartic).
\item For $E \times E'$, $\mathrm{NS} \cong \ZZ^2 \oplus \Hom(E, E')$ (Exercise~\ref{algebraic-groups:exercise-ns-product}), which by Theorem~\ref{hodge-conjecture:theorem-lefschetz-11} is also the
lattice of integral classes of type $(1,1)$, a statement about Hodge structures (Example~\ref{hodge-conjecture:example-elliptic-product}).
\end{enumerate}
\end{example}

\begin{example}\label{hodge-conjecture:example-quadric}
Let $X = \PP^1 \times \PP^1$, which is isomorphic to a nonsingular quadric surface in $\PP^3$ (Example~\ref{complex-manifolds:example-adjunction}), with projections $\mathrm{pr}_1$, $\mathrm{pr}_2$, and let $h$ be the
positive generator of $H^2(\PP^1; \ZZ)$. By the K\"unneth formula (Theorem~\ref{singular-homology:theorem-kunneth}), $H^2(X; \ZZ) = \ZZ a \oplus \ZZ b$ with $a = \mathrm{pr}_1^*h$ and $b = \mathrm{pr}_2^*h$. By
naturality of $c_1$ (Theorem~\ref{vector-bundles:theorem-line-bundles}), $a = c_1(\mathrm{pr}_1^*\mathcal{O}(1))$ and $b = c_1(\mathrm{pr}_2^*\mathcal{O}(1))$; these are the classes of the divisors
$\{p\} \times \PP^1$ and $\PP^1 \times \{p\}$, the two rulings of the quadric. So every integral class is algebraic and $\rho(X) = b_2(X) = 2$, without any Hodge theory. Conversely this determines the Hodge
numbers: $\rho(X) \leq h^{1,1}(X)$ (Corollary~\ref{hodge-conjecture:corollary-neron-severi}) and $h^{2,0} + h^{1,1} + h^{0,2} = b_2 = 2$ with $h^{2,0} = h^{0,2}$ force $h^{1,1} = 2$ and $h^{2,0} = 0$. The
intersection numbers are $a^2 = b^2 = 0$ and $ab = 1$: two lines of the same ruling are disjoint, and two lines of different rulings meet transversally in one point.
\end{example}

\begin{theorem}[Noether--Lefschetz]\label{hodge-conjecture:theorem-noether-lefschetz}
For $d \geq 4$, a very general nonsingular surface of degree $d$ in $\PP^3$ has $\Pic(X) = \ZZ\cdot\mathcal{O}_X(1)$, so $\rho(X) = 1$.
\end{theorem}

This is proved in \cite{voisin-hodge-2} and \cite{griffiths-harris-noether-lefschetz}. The proof shows, with the infinitesimal variation of Hodge structure (Chapter~\ref{vhs}), that
the locus of surfaces on which a given integral class stays of type $(1,1)$ is a proper subvariety of the parameter space, unless the class is a multiple of the hyperplane class; a
very general surface lies outside all these countably many loci. For $d \leq 3$ the statement is false: quadrics contain two families of lines, and cubic surfaces contain $27$ lines.

\section{Why the proof does not generalise}\label{hodge-conjecture:section-why-not}

\begin{remark}\label{hodge-conjecture:remark-why-not}
The proof of Theorem~\ref{hodge-conjecture:theorem-lefschetz-11} has three steps, and none generalises to codimension $p \geq 2$.
\begin{enumerate}
\item \emph{A cohomological object parametrising cycles.} Divisors modulo linear equivalence are line bundles, $\Pic(X) = H^1(X, \mathcal{O}_X^*)$, the cohomology of a sheaf of abelian
groups, and the exponential sequence relates it to $H^2(X; \ZZ)$. For codimension $p$ the Chow group $\CH^p(X) \cong H^p(X, \mathcal{K}_p)$ (Bloch's formula) involves a sheaf of Quillen
$K$-groups, which has no comparison with singular cohomology that is exact in the needed way.
\item \emph{The key lemma.} The obstruction to lifting an integral class is its image in $H^2(X, \mathcal{O}_X) = H^{0,2}$, a coherent cohomology group, and it vanishes exactly for type $(1,1)$. For
$p \geq 2$, the analogous vanishing of the components in $H^{p-1,p+1}, \ldots, H^{0,2p}$ has no sheaf-theoretic interpretation producing an object like a line bundle. Deligne cohomology
$H^{2p}_{\mathcal{D}}(X; \ZZ(p))$ fits into an exact sequence $0 \to J^p(X) \to H^{2p}_{\mathcal{D}}(X; \ZZ(p)) \to \mathrm{Hdg}^p(X; \ZZ) \to 0$, so every integral Hodge class lifts to Deligne
cohomology (Chapter~\ref{vhs}); the missing step is to lift further to cycles.
\item \emph{Sections of line bundles.} A line bundle twisted by a high multiple of an ample bundle has sections, whose zero loci are divisors. In codimension $\geq 2$ there is no
analogue: a vector bundle of rank $p$ with enough sections gives cycles as zero loci of sections, but the Chern classes of such bundles need not span the Hodge classes, and producing a
bundle with a prescribed Chern class is the vector bundle form of the conjecture itself (Proposition~\ref{hodge-conjecture:proposition-equivalent-forms}(3)).
\end{enumerate}
\end{remark}

\section{Exercises on The Lefschetz (1,1) Theorem}\label{hodge-conjecture:section-exercises-lefschetz-11}

\begin{exercise}\label{hodge-conjecture:exercise-picard-number-torus}
Let $T$ be a complex torus of dimension $2$ with $\rho(T) = 0$. Show that $T$ carries no nonconstant meromorphic functions and no curves.
\end{exercise}

\begin{solution}
A curve $C$ would give a line bundle $\mathcal{O}(C)$ with $c_1(\mathcal{O}(C)) = \mathrm{cl}(C) \neq 0$, since $\int_T\mathrm{cl}(C) \wedge \omega = \int_C\omega > 0$ for a K\"ahler form $\omega$
(Proposition~\ref{kahler:proposition-kahler-class-powers}). So $\mathrm{NS}(T) \neq 0$, a contradiction. A nonconstant meromorphic function $f$ has a nonzero polar divisor, which is a union of
curves.
\end{solution}

\begin{exercise}\label{hodge-conjecture:exercise-hodge-classes-surface}
Let $S$ be a nonsingular projective surface with $p_g(S) = h^{2,0}(S) = 0$. Show that every class in $H^2(S; \ZZ)$ is the class of a divisor, and that $\rho(S) = b_2(S)$.
\end{exercise}

\begin{solution}
If $h^{2,0} = h^{0,2} = 0$, then $H^2(S; \CC) = H^{1,1}$, so every integral class is an integral Hodge class, hence algebraic by Theorem~\ref{hodge-conjecture:theorem-lefschetz-11}. So
$\mathrm{NS}(S) = H^2(S; \ZZ)$ and $\rho = b_2$.
\end{solution}

\begin{exercise}\label{hodge-conjecture:exercise-cubic-threefold}
Show that on a nonsingular cubic threefold $X \subseteq \PP^4$ the Hodge conjecture holds, and compute $\mathrm{Hdg}^1$ and $\mathrm{Hdg}^2$.
\end{exercise}

\begin{solution}
By Corollary~\ref{hodge-conjecture:corollary-codimension-n-1} the conjecture holds in dimension $3$. By the Lefschetz hyperplane theorem and Poincar\'e duality, $H^2 = \QQ h$ and $H^4 = \QQ h^2/3$
(the class of a line, since $\int_Xh^3 = 3$), both algebraic, so $\mathrm{Hdg}^1 = H^2(X; \QQ)$ and $\mathrm{Hdg}^2 = H^4(X; \QQ)$. The interesting group $H^3$, of dimension $10$ with $h^{2,1} = 5$, has
no Hodge classes, being of odd degree.
\end{solution}
